% math.tex — algorithms module
% Mathematical results supporting crates/src/algorithms/
%
% Labels defined here:
%   alg:ehz, alg:billiard, alg:tube,
%   cor:adjacency-pruning,
%   def:tube, def:tube-data, def:tube-extension, def:tube-close,
%   def:rotation-increment, def:symplectic-polytope,
%   lem:base-point-recovery, lem:fixed-point, lem:lagrangian-facets,
%   lem:prune-action, lem:prune-empty, lem:prune-rotation, lem:prune-simple,
%   lem:shoelace (cross-ref to geom/math.tex),
%   lem:sigma-structure,
%   rem:beta-to-tau,
%   thm:billiard-characterization, thm:bounce-bound,
%   thm:conformality, thm:sympl-invariance
%
% Labels referenced but defined elsewhere (do NOT redefine):
%   lem:kkt (kkt/math.tex), def:ehz-capacity (geom/math.tex),
%   thm:hko-counterexample (geom/math.tex),
%   lem:numerical-transition-feasibility (kkt/math.tex),
%   lem:q-error-bound (kkt/math.tex)
%
% Sources: thesis/general-case-algorithm.tex,
%   thesis/pruned-general-case-algorithm.tex,
%   thesis/lagrangian-product-algorithm.tex,
%   thesis/lagrangian-product-algorithm-proof.tex,
%   thesis/tube-algorithm.tex,
%   thesis/simple-minimizer-existence.tex,
%   .rs doc comments (stubs for labels not yet in thesis)

% \input'd from crates/src/math.tex — do not add \documentclass or \begin{document}

% ====================================================================
% Part 1: General-case algorithm (HK2017)
% ====================================================================

\begin{algorithm}[EHZ capacity]\label{alg:ehz}
% Jörn: text approved (06a976c)
% Source: thesis/general-case-algorithm.tex
% Parameterization: K = {x : a_i^T x ≤ 1}, dual vertices a_i
Computes $c_{\mathrm{EHZ}}(K)$ via exhaustive enumeration.

\smallskip\noindent
\emph{Input:} $a_1, \ldots, a_F$ (dual vertices).
\quad
\emph{Output:} $c_{\mathrm{EHZ}}(K)$.

\smallskip\noindent
For each nonempty subset $S \subseteq \{1, \ldots, F\}$
and each permutation $\sigma$ of $S$
(identified up to cyclic shifts):

\begin{enumerate}
\item \textbf{Define}
  the dual-vertex matrix $A \in \mathbb{R}^{|S| \times 4}$
  and action matrix $H \in \mathbb{R}^{|S| \times |S|}$:
  \begin{align*}
    A_{id} &= a_{\sigma(i),d},
      & i &= 1, \ldots, |S|, \quad d = 1, \ldots, 4, \\
    H_{ij} &= \begin{cases}
      \omega_0(a_{\sigma(i)}, a_{\sigma(j)}) & \text{if } i < j, \\
      0 & \text{if } i = j, \\
      \omega_0(a_{\sigma(j)}, a_{\sigma(i)}) & \text{if } i > j.
    \end{cases}
  \end{align*}

\item \textbf{Solve.}
  \[
    \begin{pmatrix} H & A & \mathbf{1} \\ A^\top & 0 & 0 \\ \mathbf{1}^\top & 0 & 0 \end{pmatrix}
    \begin{pmatrix} \beta \\ \mu \\ \xi \end{pmatrix}
    = \begin{pmatrix} 0 \\ 0 \\ 1 \end{pmatrix}.
  \]

\item \textbf{Filter.}
  Discard $(S, \sigma)$ if the system has no solution, or
  if $\beta_i \leq 0$ for any~$i$.
  If the system has multiple solutions, choose any solution
  with $\beta_i > 0$ for all~$i$, if one exists.

\item \textbf{Evaluate.}
  Record $\beta^\top H \beta$.
\end{enumerate}

\noindent
Return
$c_{\mathrm{EHZ}}(K) = \bigl(\max_{S,\sigma} \beta^\top H \beta\bigr)^{-1}$.
\end{algorithm}


\begin{corollary}[Adjacency pruning]\label{cor:adjacency-pruning}
% Jörn: text approved (0839595)
% Source: thesis/pruned-general-case-algorithm.tex
In Algorithm~\ref{alg:ehz}, discard any
$(S, \sigma)$ for which some consecutive pair
$\sigma(k), \sigma(k{+}1)$ (cyclically) satisfies
$F_{\sigma(k)} \cap F_{\sigma(k+1)} = \emptyset$.
Equivalently, restrict to cyclic orderings~$\sigma$ of~$S$ that
form a Hamiltonian cycle in the induced subgraph~$G_K[S]$.
\end{corollary}


\begin{unverified}
\begin{theorem}[Conformality]\label{thm:conformality}
% [TODO: JÖRN - verify statement, add proof]
% Source: hk2017/conformality_test.rs doc comment
EHZ capacity is degree-2 homogeneous: for any convex body $K$
and $\alpha > 0$,
\[
  c_{\mathrm{EHZ}}(\alpha K) = \alpha^2\, c_{\mathrm{EHZ}}(K).
\]
\end{theorem}
\end{unverified}


\begin{unverified}
\begin{theorem}[Symplectic invariance]\label{thm:sympl-invariance}
% [TODO: JÖRN - verify statement, add proof]
% Source: hk2017/symplectic_invariance_test.rs doc comment
EHZ capacity is invariant under symplectomorphisms: for any
$M \in \operatorname{Sp}(4)$ and $b \in \mathbb{R}^4$,
\[
  c_{\mathrm{EHZ}}(MK + b) = c_{\mathrm{EHZ}}(K).
\]
\end{theorem}
\end{unverified}


% ====================================================================
% Part 2: Orbit recovery
% ====================================================================

\begin{lemma}[Base point recovery]\label{lem:base-point-recovery}
% Jörn: math approved (f7d26ea)
% Source: thesis/simple-minimizer-existence.tex
% Parameterization: K = {x : a_i^T x ≤ 1}, dual vertices a_i
Let $K = \{x : \langle a_i, x \rangle \leq 1\}$ be a polytope
and $(\sigma, \tau, b)$ a simple orbit
with $m$ active facets
$S = \{\sigma(k) : \tau_k > 0\}$.
Define the accumulated displacements
\[
  v_0 = 0, \qquad
  v_k = \sum_{j=1}^{k-1} \tau_j\, R_{\sigma(j)}
  \quad \text{for } k = 1, \ldots, m.
\]
A point $b \in \mathbb{R}^4$ is a valid base point
(i.e., the orbit $\gamma$ with $\gamma(0) = b$ and
velocity $R_{\sigma(k)}$ on the $k$-th segment
is a generalized Reeb orbit on~$\partial K$) if and only if:
\begin{enumerate}
\item \emph{On-facet (equality):}
  for each active $k$,
  $\langle a_{\sigma(k)},\, b \rangle
    = 1 - \langle a_{\sigma(k)},\, v_k \rangle$;
\item \emph{Inside $K$ (inequality):}
  for each active segment $k$ and all $s \in [0, \tau_k]$,
  $\langle a_j,\, b + v_k + s\, R_{\sigma(k)} \rangle
    \leq 1$ for all $j = 1, \ldots, F$.
\end{enumerate}
The equality constraints form
a linear system $A_S\, b = r$ with $m$ equations and $4$~unknowns.
The set of valid base points is a convex polytope of
dimension at most $4 - \operatorname{rank}(A_S)$.
\end{lemma}

\begin{proof}
\emph{Equality constraints.}
During the $k$-th segment, the orbit position is
\[
  \gamma(t) = b + v_k + (t - t_k)\, R_{\sigma(k)},
  \quad t \in [t_k,\, t_k + \tau_k],
\]
where $t_k = \sum_{j=1}^{k-1} \tau_j$.
For $\gamma$ to lie on facet $F_{\sigma(k)}$,
we require $\langle a_{\sigma(k)},\, \gamma(t) \rangle
= 1$ throughout the segment.
Since $R_{\sigma(k)}$ is tangent to $F_{\sigma(k)}$
($\langle a_{\sigma(k)}, R_{\sigma(k)} \rangle = 0$),
the time-dependent term vanishes:
\[
  \langle a_{\sigma(k)},\, \gamma(t) \rangle
  = \langle a_{\sigma(k)},\, b + v_k \rangle.
\]
Hence the on-facet condition for the entire segment
reduces to a single equation.

\emph{Inequality constraints.}
The orbit must remain inside~$K$:
$\langle a_j, \gamma(t) \rangle \leq 1$
for all $j$ and all~$t$.
Since $\gamma(t)$ is affine in~$t$ on each segment,
each inequality is linear in~$s$.
A linear function on $[0, \tau_k]$ attains its
maximum at an endpoint, so it suffices to
check at $s = 0$ and $s = \tau_k$ only.

\emph{Solution structure.}
The equality constraints are a linear system
$A_S\, b = r$ in~$b \in \mathbb{R}^4$.
Its solution set is an affine subspace of
dimension $4 - \operatorname{rank}(A_S)$.
The inequality constraints are linear in~$b$,
so the intersection is a convex polytope.
\end{proof}


\begin{remark}[Conversion from algorithm output]\label{rem:beta-to-tau}
% Jörn: math approved (f7d26ea)
% Source: thesis/simple-minimizer-existence.tex
% Parameterization: K = {x : a_i^T x ≤ 1}, dual vertices a_i
Algorithm~\ref{alg:ehz} outputs a permutation~$\sigma$ and
a vector~$\beta$ with $\beta_i \geq 0$ and
$\sum_i \beta_i = 1$.
The dwell times are
\[
  \tau_k = T\, \beta_{\sigma(k)},
\]
where $T = c_{\mathrm{EHZ}}(K)$ is the capacity
(the period of the orbit).
\end{remark}


% ====================================================================
% Part 3: Billiard algorithm (Lagrangian products)
% ====================================================================

\begin{unverified}
\begin{lemma}[Facet structure of Lagrangian products]\label{lem:lagrangian-facets}
% Source: thesis/lagrangian-product-algorithm-proof.tex
Let $K = K_q \times K_p$ be a Lagrangian product.
Let $K_q = \{x \in L_q : \langle a_k^q, x \rangle \leq 1\}$
with dual vertices
$a_1^q, \ldots, a_{F_q}^q \in L_q$,
and similarly $K_p$ with dual vertices
$a_1^p, \ldots, a_{F_p}^p \in L_p$.
Then $K$ has $F = F_q + F_p$ facets, partitioned into:
\begin{itemize}
\item \emph{$q$-facets} with dual vertices $a_i = (a_i^q, 0)$,
  for $i = 1, \ldots, F_q$;
\item \emph{$p$-facets} with dual vertices $a_j = (0, a_j^p)$,
  for $j = 1, \ldots, F_p$.
\end{itemize}
Every dual vertex lies entirely in $L_q$ or entirely in $L_p$.
\end{lemma}

\begin{proof}
A facet of $K_q \times K_p$ arises from a facet (edge) of
one factor and the full other factor: either
$\text{edge}_i(K_q) \times K_p$ or $K_q \times \text{edge}_j(K_p)$.
The defining constraint $\langle a_i^q, x_q \rangle \leq 1$
for edge~$i$ of~$K_q$ lifts to
$\langle (a_i^q, 0),\, x \rangle \leq 1$,
so the dual vertex is $(a_i^q, 0)$;
similarly for $p$-facets.
\end{proof}
\end{unverified}


\begin{unverified}
\begin{lemma}[Sigma structure for Lagrangian products]\label{lem:sigma-structure}
% Source: thesis/lagrangian-product-algorithm-proof.tex
Let $K = K_q \times K_p$ be a Lagrangian product, and let
$(\sigma, \tau)$ be a simple Reeb orbit on $\partial K$
with all $\tau_k > 0$.
Write the cyclic sequence $\sigma$ as a word in the alphabet
$\{Q, P\}$ by replacing each $q$-facet index with $Q$ and each
$p$-facet index with $P$.
Then:
\begin{enumerate}
\item[(i)] The word has the form
  $\bigl([Q \mid QQ]\,[P \mid PP]\bigr)^k$ for some $k \geq 1$:
  alternating blocks of $q$- and $p$-facets, each block of length
  one or two.

\item[(ii)] Within each block of length two, the two facets are
  adjacent in $G_K$
  (consecutive edges of $K_q$ or $K_p$ sharing a vertex).
\end{enumerate}
\end{lemma}

\begin{proof}
\textbf{Claim: no three consecutive same-type facets.}\quad
Suppose for contradiction that $\sigma$ contains three
consecutive $q$-facets $q_1, q_2, q_3$ (all with $\tau > 0$).
The Reeb vectors $R_{q_1}, R_{q_2}, R_{q_3}$ all lie in~$L_p$,
so $q$ is constant throughout the three segments.
The orbit at the breakpoint between $q_1$ and $q_2$ lies on
both facets, meaning $\langle a_{q_1}^q, q \rangle = 1$
and $\langle a_{q_2}^q, q \rangle = 1$: the point $q$
lies on both edges $1$ and $2$ of~$K_q$, hence at their
common vertex $v_{12}$.
Similarly, the breakpoint between $q_2$ and $q_3$ forces $q$
to lie at vertex $v_{23}$ (intersection of edges $2$ and $3$).
Since $q$ is constant, $v_{12} = v_{23}$.
But for a nondegenerate convex polygon, $v_{12}$ and $v_{23}$
are the two distinct endpoints of edge~$2$: contradiction.

The same argument applies to $p$-facets by symmetry.

\medskip\noindent
\textbf{Alternation and block structure.}\quad
Since no three consecutive facets have the same type, each maximal
run of same-type facets has length one or two.
% [TODO: JÖRN - alternation argument is imprecise: two adjacent length-1
%   q-runs would give only 2 consecutive q-facets, not 3. The correct
%   reason is that adjacent same-type maximal runs would merge.]
Consecutive runs must alternate between $q$ and $p$.
In a cyclic sequence, alternation forces equal block counts:
$k$ blocks of $q$-facets and $k$ blocks of $p$-facets.

\medskip\noindent
\textbf{Within-block adjacency.}\quad
Follows directly from the adjacency constraint:
consecutive facets in $\sigma$ share a breakpoint, hence are
adjacent in~$G_K$.
Two same-type adjacent facets correspond to edges sharing a
vertex in the 2D polygon.
\end{proof}
\end{unverified}


\begin{unverified}
\begin{theorem}[Billiard characterization of EHZ capacity]%
\label{thm:billiard-characterization}
% Source: thesis/lagrangian-product-algorithm-proof.tex
% Attribution: Artstein-Avidan \& Ostrover (2014); Rudolf (2022).
For a Lagrangian product $K = K_q \times K_p$:
\[
  c_{\mathrm{EHZ}}(K_q \times K_p)
  = \min\bigl\{ \ell_{K_p^\circ}(\eta) :
  \eta \text{ is a closed $K_p$-billiard trajectory in } K_q \bigr\},
\]
where $\ell_{K_p^\circ}(\eta) = \sum_i h_{K_p}(\Delta_i)$ is the
$K_p^\circ$-length and $\Delta_i = x_{i+1} - x_i$ are the edge vectors.
\end{theorem}
% [TODO: JÖRN - proof is in the referenced papers, not reproduced here]
\end{unverified}


\begin{unverified}
\begin{theorem}[Bounce bound]%
\label{thm:bounce-bound}
% Source: thesis/lagrangian-product-algorithm-proof.tex
% Attribution: Bezdek \& Bezdek (2009); Rudolf (2022).
For convex polygons $K_q, K_p \subset \mathbb{R}^2$,
the minimum in Theorem~\ref{thm:billiard-characterization}
is achieved by a billiard trajectory with at most $3$ bounce points.
\end{theorem}
% [TODO: JÖRN - proof is in the referenced papers, not reproduced here]
\end{unverified}


\begin{unverified}
\begin{algorithm}[Billiard capacity]\label{alg:billiard}
% Source: thesis/lagrangian-product-algorithm.tex
Computes $c_{\mathrm{EHZ}}(K_q \times K_p)$
via Theorem~\ref{thm:billiard-characterization}.

\emph{Input:} Lagrangian product $K = K_q \times K_p$ in
$H$-representation.
\quad
\emph{Output:} $c_{\mathrm{EHZ}}(K)$.

\begin{enumerate}
\item \textbf{Classify} facets into $q$-type and $p$-type
  (Lemma~\ref{lem:lagrangian-facets}).
\item \textbf{Build blocks.}
  For each polygon ($K_q$ and $K_p$), enumerate all
  \emph{single-facet blocks} (one edge) and
  \emph{pair blocks} (two adjacent edges, both orderings).
% [TODO: JÖRN - k=1 is excluded here. The thesis proves this via
%   lem:k1-infeasible (closure condition forces antiparallel normals
%   on adjacent edges, impossible for convex polygons). Add citation or
%   copy the lemma.]
\item \textbf{Enumerate.}
  For $k = 2, 3$: for each selection of $k$ non-overlapping
  $q$-blocks and $k$ non-overlapping $p$-blocks, and each
  interleaving into the pattern
  $\bigl([Q \mid QQ]\,[P \mid PP]\bigr)^k$,
  flatten the blocks into a facet sequence~$\sigma$.
\item \textbf{Solve and filter.}
  For each $\sigma$: solve the KKT system,
  filter by $\beta_i > 0$ and $Q(\beta) > 0$, compute
  $A(S, \sigma) = \tfrac{1}{2}\, Q(\beta)^{-1}$.
\item \textbf{Return} $\min A(S, \sigma)$.
\end{enumerate}
\end{algorithm}
\end{unverified}


% ====================================================================
% Part 4: Tube algorithm
% ====================================================================

\begin{unverified}
\begin{definition}[Symplectic polytope]\label{def:symplectic-polytope}
% Source: thesis/tube-algorithm.tex
% Parameterization: K = {x : a_i^T x ≤ 1}, dual vertices a_i
A \emph{Lagrangian $2$-face} is a $2$-face $F_{ij}$ on which
$\omega_0|_{T F_{ij}} = 0$, equivalently $\omega_0(a_i, a_j) = 0$.
A \emph{symplectic polytope} is a polytope $K \subset \mathbb{R}^4$
that has no Lagrangian $2$-face:
\[
  \omega_0(a_i, a_j) \neq 0
  \qquad \text{for all pairs } (i,j)
  \text{ such that } F_i \cap F_j \text{ is a $2$-face.}
\]
\end{definition}
\end{unverified}


\begin{unverified}
\begin{definition}[Tube]\label{def:tube}
% Source: thesis/tube-algorithm.tex
The \emph{tube} $T(\sigma)$ of a type $\sigma = (\sigma(1), \ldots, \sigma(k))$
is the set of all pure Reeb trajectories of type $\sigma$:
\[
  T(\sigma) \coloneqq \bigl\{
    \gamma \colon \mathbb{R} \to \partial K
    : \gamma \text{ is a pure Reeb trajectory of type } \sigma
  \bigr\}.
\]
A pure Reeb trajectory of type $\sigma$ has
$\dot{\gamma}(t) = R_{\sigma(i)}$ on each interval
$(t_{i-1}, t_i)$ for $i = 1, \ldots, k$,
where $F_{\sigma(i)} \to F_{\sigma(i+1)}$ is a directed $2$-face
for each $i = 1, \ldots, k-1$.
\end{definition}
\end{unverified}


\begin{unverified}
\begin{definition}[Tube data structure]\label{def:tube-data}
% Source: thesis/tube-algorithm.tex
A tube $T(\sigma)$ is represented by:
\begin{enumerate}
\item $\operatorname{Start} \subset F_{\sigma(1)} \cap F_{\sigma(2)}$:
  closed convex polygon.
\item $\operatorname{End} \subset F_{\sigma(k-1)} \cap F_{\sigma(k)}$:
  closed convex polygon.
\item $\varphi \colon \operatorname{Start} \to \operatorname{End}$:
  the affine map sending $\gamma(t_1)$ to $\gamma(t_{k-1})$.
\item $a \colon \operatorname{End} \to \mathbb{R}$:
  the affine function giving the total elapsed time (= action)
  of $\gamma|_{[t_1, t_{k-1}]}$.
\item $\rho \in \mathbb{R}$:
  the rotation number accumulated on $(t_1, t_{k-1})$.
\end{enumerate}
\end{definition}
\end{unverified}


\begin{unverified}
\begin{definition}[Tube extension]\label{def:tube-extension}
% Source: thesis/tube-algorithm.tex
Let $T_k = T(\sigma(1), \ldots, \sigma(k))$ be a tube
and $\sigma(k+1)$ satisfy $F_{\sigma(k)} \to F_{\sigma(k+1)}$.
Write $\Phi \coloneqq \Phi_{\sigma(k-1),\sigma(k),\sigma(k+1)}$
for the step map.
The extended tube $T(\sigma(1), \ldots, \sigma(k+1))$ has data:
\begin{align*}
  \operatorname{End}' &\coloneqq
    \Phi(\operatorname{End}) \cap \bigl( F_{\sigma(k)} \cap F_{\sigma(k+1)} \bigr), \\
  \varphi' &\coloneqq \Phi \circ \varphi, \\
  \operatorname{Start}' &\coloneqq (\varphi')^{-1}(\operatorname{End}'), \\
  a'(y') &\coloneqq
    a\bigl(\Phi^{-1}(y')\bigr)
    + \Delta a_{\sigma(k-1),\sigma(k),\sigma(k+1)}\bigl(\Phi^{-1}(y')\bigr), \\
  \rho' &\coloneqq \rho + \Delta\rho_{\sigma(k),\sigma(k+1)}.
\end{align*}
\end{definition}
\end{unverified}


\begin{unverified}
\begin{definition}[Rotation increment at a breakpoint]\label{def:rotation-increment}
% Source: thesis/tube-algorithm.tex
Let $F_j \to F_l$ be a directed $2$-face.
The \emph{rotation increment} $\Delta\rho_{jl} \in \mathbb{R}$
is the contribution to the rotation number at the breakpoint
where a Reeb trajectory transitions from velocity $R_j$ to $R_l$.
It is computed from the \emph{transition matrix}
$\psi_{F_{jl}} \in \operatorname{Sp}(2)$.
By the theory of CH2021, $\psi_{F_{jl}}$ is positive elliptic
and $\Delta\rho_{jl} \in (0, \tfrac{1}{2})$.
\end{definition}
\end{unverified}


\begin{unverified}
\begin{definition}[Closing a tube]\label{def:tube-close}
% Source: thesis/tube-algorithm.tex
Given a tube $T(\sigma(1), \ldots, \sigma(k))$,
its \emph{closure} is obtained by performing two extension steps:
appending $\sigma(k+1) = \sigma(1)$ and then $\sigma(k+2) = \sigma(2)$.
The result is a tube $T(\sigma(1), \ldots, \sigma(k), \sigma(1), \sigma(2))$
with flow map $\varphi'' \colon \operatorname{Start} \to F_{\sigma(1)} \cap F_{\sigma(2)}$
and action function $a'' \colon \operatorname{End}'' \to \mathbb{R}$.
\end{definition}
\end{unverified}


% ── Pruning lemmas ──

\begin{unverified}
\begin{lemma}[Pruning: empty tube]\label{lem:prune-empty}
% Source: thesis/tube-algorithm.tex
If $\operatorname{Start}(\sigma) = \emptyset$ or $\operatorname{End}(\sigma) = \emptyset$,
then $T(\sigma) = \emptyset$.
\end{lemma}

\begin{proof}
A trajectory $\gamma \in T(\sigma)$ satisfies
$\gamma(t_1) \in \operatorname{Start}(\sigma)$ and
$\gamma(t_{k-1}) \in \operatorname{End}(\sigma)$ by definition.
\end{proof}
\end{unverified}


\begin{unverified}
\begin{lemma}[Pruning: action lower bound]\label{lem:prune-action}
% Source: thesis/tube-algorithm.tex
Let $c^* > 0$ be the action of the current best candidate.
If $\min_{y \in \operatorname{End}} a(y) > c^*$,
then $T(\sigma)$ contains no minimum-action orbit.
\end{lemma}

\begin{proof}
% [TODO: JÖRN - the inequality A(gamma) >= a(gamma(t_{k-1})) asserts that
%   closing a tube adds non-negative action. This is plausible but not proved here.]
Every $\gamma \in T(\sigma)$ has
$A(\gamma) \geq a(\gamma(t_{k-1})) \geq \min_{\operatorname{End}} a > c^*$.
Since the minimum-action orbit has action $\leq c^*$, it is not in $T(\sigma)$.
\end{proof}
\end{unverified}


\begin{unverified}
\begin{lemma}[Pruning: rotation upper bound]\label{lem:prune-rotation}
% Source: thesis/tube-algorithm.tex
Let $\rho$ be the rotation number of tube $T(\sigma)$.
If $\rho \geq 2$, then $T(\sigma)$ contains no minimum-action orbit.
\end{lemma}

\begin{proof}
By CH2021, the minimum-action combinatorial Reeb orbit satisfies
$\rho_{\mathrm{comb}} \leq 2$.
The closing step adds at least two
rotation increments, each in $(0, \tfrac{1}{2})$,
so any closed orbit from $T(\sigma)$ has total rotation
$> \rho \geq 2$.
Hence no closed orbit from $T(\sigma)$ satisfies
$\rho_{\mathrm{comb}} \leq 2$.
\end{proof}
\end{unverified}


\begin{unverified}
\begin{lemma}[Pruning: simple orbit]\label{lem:prune-simple}
% Source: thesis/tube-algorithm.tex
If $\sigma(i) = \sigma(j)$ for some $1 \leq i < j \leq k$,
then $T(\sigma)$ contains no minimum-action orbit.
\end{lemma}

\begin{proof}
% [TODO: JÖRN - cite HK2017 or CH2021 for simplicity of the minimum-action orbit.]
The minimum-action generalized Reeb orbit is simple:
it visits each facet on a connected time interval and does not repeat facets.
A pure Reeb trajectory of type $\sigma$, when closed,
visits $F_{\sigma(i)}$ twice if $\sigma(i) = \sigma(j)$.
Hence no closed orbit from $T(\sigma)$ can be the minimum-action orbit.
\end{proof}
\end{unverified}


\begin{unverified}
\begin{lemma}[Closed orbits as fixed points]\label{lem:fixed-point}
% Source: thesis/tube-algorithm.tex
The closed Reeb orbits in the closure of $T(\sigma)$ are in bijection with
fixed points of $\varphi'' \colon \operatorname{Start}' \to F_{\sigma(1)} \cap F_{\sigma(2)}$:
\[
  \gamma \text{ is a closed orbit in } T(\sigma)
  \;\Longleftrightarrow\;
  \varphi''(\gamma(t_1)) = \gamma(t_1)
  \;\text{ and }\;
  \gamma(t_1) \in \operatorname{Start}'.
\]
The action of the closed orbit through the fixed point $x$ is $a''(x)$.
\end{lemma}

\begin{proof}
% [TODO: JÖRN - check time index: def:tube-close appends two facets
%   (sigma(k+1) and sigma(k+2)), so the final time may be t_{k+2}, not t_{k+1}.]
A pure Reeb trajectory $\gamma \in T(\sigma(1), \ldots, \sigma(k))$
is closed iff $\gamma(t_{k+1}) = \gamma(t_1)$.
After closure,
\[
  \gamma(t_{k+1}) = \varphi''(\gamma(t_1)),
\]
so the periodicity condition is $\varphi''(\gamma(t_1)) = \gamma(t_1)$.
The action is $a''(\varphi''(\gamma(t_1))) = a''(\gamma(t_1))$.
\end{proof}
\end{unverified}


\begin{unverified}
\begin{algorithm}[Tube algorithm]\label{alg:tube}
% Source: thesis/tube-algorithm.tex
Computes $c_{\mathrm{EHZ}}(K)$ for a symplectic polytope $K$
(assuming no Type~2 orbits).

\smallskip\noindent
% Parameterization: K = {x : a_i^T x ≤ 1}, dual vertices a_i
\emph{Input:} Dual vertices $a_1, \ldots, a_F$.
\quad
\emph{Output:} $c_{\mathrm{EHZ}}(K)$.

\smallskip\noindent
\textbf{Precomputation.}
\begin{enumerate}
\item For each pair $(i,j)$: check if $F_i \cap F_j$ is a $2$-face
  and compute $\omega_0(a_i, a_j)$. Record directed edges $F_i \to F_j$
  where $\omega_0(a_i, a_j) > 0$.
\item For each valid triple $(i,j,l)$ with $F_i \to F_j$ and $F_j \to F_l$:
  compute the step map $\Phi_{ijl}$,
  action increment $\Delta a_{ijl}$,
  and rotation increment $\Delta\rho_{jl}$.
\end{enumerate}

\smallskip\noindent
\textbf{Search.}
Initialize $c^* = +\infty$.
For each directed edge $(i, j)$ with $F_i \to F_j$,
initialize tube $T(i, j)$.
Then recurse: given a tube $T(\sigma(1), \ldots, \sigma(k))$:

\begin{enumerate}
\item \textbf{Prune.} Discard if any of the following hold:
  \begin{itemize}
    \item $\operatorname{Start} = \emptyset$ or $\operatorname{End} = \emptyset$
      (Lemma~\ref{lem:prune-empty});
    \item $\min_{y \in \operatorname{End}} a(y) > c^*$
      (Lemma~\ref{lem:prune-action});
    \item $\rho \geq 2$ (Lemma~\ref{lem:prune-rotation}).
  \end{itemize}

\item \textbf{Try closing.}
  Form the closure
  (Definition~\ref{def:tube-close}).
  Find fixed points $x$ of $\varphi''$ in $\operatorname{Start}'$
  (Lemma~\ref{lem:fixed-point}).
  For each fixed point $x$, update $c^* \leftarrow \min(c^*, a''(x))$.

\item \textbf{Extend.}
  For each $l \notin \{\sigma(1), \ldots, \sigma(k)\}$
  with $F_{\sigma(k)} \to F_l$
  (Lemma~\ref{lem:prune-simple}):
  compute $T(\sigma(1), \ldots, \sigma(k), l)$
  (Definition~\ref{def:tube-extension}) and recurse.
\end{enumerate}

\noindent
Return $c^*$.
\end{algorithm}
\end{unverified}

